%========================================================
% APPENDIX B.1 — Basis Dependence, Representation Dependence,
% and Invariance of Concurrence
%========================================================

\section{Concurrence: Basis Dependence, Representation Dependence, and Invariance}

This appendix clarifies the distinction between representation-dependent quantities and physically invariant quantities in quantum mechanics, with particular emphasis on the concurrence used throughout the present work.

\medskip

A potential conceptual ambiguity arises from the standard definition of concurrence for mixed two-qubit states, since the construction explicitly involves the operator \( \sigma_y \) and complex conjugation in the computational basis. At first sight, this may suggest that concurrence depends on a preferred polarization axis or on a specific basis representation.

The purpose of this appendix is to show that this interpretation is incorrect, and that concurrence constitutes an intrinsic, basis-independent measure of bipartite entanglement.

%--------------------------------------------------------
% D.1 - Representation Dependence and Basis Transformations
%--------------------------------------------------------

\subsection*{Representation Dependence and Basis Transformations}

In quantum mechanics, vectors and operators are represented through coordinates and matrices relative to a chosen basis.

For a single qubit, a generic pure state may be written as

\[
|\psi\rangle
=
\alpha |0\rangle
+
\beta |1\rangle.
\]

The coefficients \( \alpha \) and \( \beta \) depend on the chosen basis representation.

For example, switching from the computational basis

\[
\{|0\rangle, |1\rangle\}
\]

to the rotated basis

\[
|+\rangle
=
\frac{|0\rangle+|1\rangle}{\sqrt{2}},
\qquad
|-\rangle
=
\frac{|0\rangle-|1\rangle}{\sqrt{2}},
\]

changes the numerical coordinates of the state vector, while leaving the physical quantum state unchanged.

\medskip

Similarly, a density operator transforms under a basis change according to

\begin{equation}
\rho
\rightarrow
\rho'
=
U \rho U^\dagger,
\label{eq:density_basis_transformation}
\end{equation}

where \( U \) is a unitary change-of-basis operator.

The individual matrix entries of \( \rho \) therefore depend on the chosen representation.

However, physically meaningful invariant quantities such as

\[
\mathrm{Tr}(\rho),
\qquad
\mathrm{Tr}(\rho^2),
\qquad
\det(\rho),
\qquad
\text{spec}(\rho)
\]

remain unchanged under Eq.~\eqref{eq:density_basis_transformation}.

\medskip

This distinction is fundamental:

\begin{itemize}

\item matrix elements and coordinate representations are generally basis dependent,

\item intrinsic physical properties are basis independent.

\end{itemize}

%--------------------------------------------------------
% D.2 - Direction Dependence and Physical Observables
%--------------------------------------------------------

\subsection*{Direction Dependence and Physical Observables}

A separate concept is the dependence of observables on physical measurement directions.

Expectation values such as

\[
\langle \sigma_x \rangle,
\qquad
\langle \sigma_y \rangle,
\qquad
\langle \sigma_z \rangle
\]

correspond to measurements performed along different axes of the Bloch sphere.

These quantities are physically direction dependent, since changing the measurement axis changes the measured observable.

\medskip

For instance,

\[
\langle \sigma_z \rangle
=
\mathrm{Tr}(\rho \sigma_z)
\]

measures polarization along the computational \( Z \)-axis, while

\[
\langle \sigma_x \rangle
=
\mathrm{Tr}(\rho \sigma_x)
\]

measures polarization along the transverse \( X \)-axis.

\medskip

Direction dependence therefore refers to the physical choice of measurement axis, rather than to the mathematical choice of representation basis.

These two notions are related but conceptually distinct.

%--------------------------------------------------------
% D.3 - Representation Dependence of Pauli Matrices
%--------------------------------------------------------

\subsection*{Representation Dependence of Pauli Matrices}

The Pauli operators themselves admit different matrix representations depending on the chosen basis.

In the computational basis,

\begin{equation}
\sigma_y
=
\begin{pmatrix}
0 & -i \\
i & 0
\end{pmatrix}.
\label{eq:sigma_y_computational}
\end{equation}

Under a basis transformation \( U \), the matrix representation changes according to

\begin{equation}
\sigma_y
\rightarrow
\sigma_y'
=
U \sigma_y U^\dagger.
\label{eq:sigma_y_basis_change}
\end{equation}

Consequently, the matrix entries appearing in Eq.~\eqref{eq:sigma_y_computational} are representation dependent.

\medskip

Importantly, this does not imply that the physical observable itself changes. Rather, only its coordinate representation changes.

This situation is directly analogous to ordinary vector coordinates in Euclidean geometry.

%--------------------------------------------------------
% D.4 - Concurrence and the Spin-Flip Construction
%--------------------------------------------------------

\subsection*{Concurrence and the Spin-Flip Construction}

For mixed two-qubit states, concurrence is commonly expressed through the spin-flipped density operator

\begin{equation}
\tilde{\rho}
=
(\sigma_y \otimes \sigma_y)
\rho^*
(\sigma_y \otimes \sigma_y),
\label{eq:appendix_spin_flip_definition}
\end{equation}

where \( \rho^* \) denotes complex conjugation in the computational basis.

The concurrence is then defined as

\begin{equation}
C(\rho)
=
\max
\left(
0,
\lambda_1-\lambda_2-\lambda_3-\lambda_4
\right),
\label{eq:appendix_concurrence_definition}
\end{equation}

where \( \lambda_i \) are the ordered square roots of the eigenvalues of

\[
\rho \tilde{\rho}.
\]

\medskip

At first sight, Eq.~\eqref{eq:appendix_spin_flip_definition} may appear to privilege the \( Y \)-direction of the Bloch sphere because of the explicit presence of \( \sigma_y \).

However, this interpretation is incorrect.

In this context, \( \sigma_y \) is not used as a measurement operator associated with the physical \( Y \)-axis. Instead, it appears because it realizes the invariant antisymmetric structure of two-dimensional spinor spaces.

%--------------------------------------------------------
% D.5 - Antisymmetric Structure of Two-Qubit States
%--------------------------------------------------------

\subsection*{Antisymmetric Structure of Two-Qubit States}

Consider a generic pure two-qubit state

\begin{equation}
|\psi\rangle
=
a|00\rangle
+
b|01\rangle
+
c|10\rangle
+
d|11\rangle.
\label{eq:generic_two_qubit_state_appendix}
\end{equation}

The coefficients may be organized into the matrix

\begin{equation}
\Psi
=
\begin{pmatrix}
a & b \\
c & d
\end{pmatrix}.
\label{eq:coefficient_matrix_appendix}
\end{equation}

For pure states, concurrence reduces to

\begin{equation}
C(|\psi\rangle)
=
2|ad-bc|.
\label{eq:pure_state_concurrence_appendix}
\end{equation}

The quantity

\[
ad-bc
\]

is precisely the determinant of the coefficient matrix \( \Psi \):

\begin{equation}
\det(\Psi)
=
ad-bc.
\label{eq:determinant_coefficient_matrix}
\end{equation}

This determinant originates from the antisymmetric bilinear form

\begin{equation}
\varepsilon
=
\begin{pmatrix}
0 & 1 \\
-1 & 0
\end{pmatrix},
\label{eq:antisymmetric_tensor_appendix}
\end{equation}

which satisfies

\begin{equation}
\sigma_y
=
-i\varepsilon.
\label{eq:sigma_y_epsilon_relation}
\end{equation}

Thus, the appearance of \( \sigma_y \) in the concurrence construction is not associated with a preferred physical polarization axis, but with the antisymmetric invariant structure of two-level quantum systems.

%--------------------------------------------------------
% D.6 - Basis Independence of Concurrence
%--------------------------------------------------------

\subsection*{Basis Independence of Concurrence}

An entanglement measure must remain invariant under local unitary transformations, since local polarization rotations cannot create or destroy entanglement.

For any local unitary evolution,

\begin{equation}
\rho
\rightarrow
(U_A \otimes U_B)
\rho
(U_A^\dagger \otimes U_B^\dagger),
\label{eq:local_unitary_appendix}
\end{equation}

the concurrence satisfies

\begin{equation}
C
\left[
(U_A \otimes U_B)
\rho
(U_A^\dagger \otimes U_B^\dagger)
\right]
=
C(\rho).
\label{eq:concurrence_local_unitary_invariance}
\end{equation}

This invariance is essential for the physical interpretation of concurrence as an intrinsic measure of bipartite entanglement.

\medskip

Consequently:

\begin{itemize}

\item the matrix representation of the spin-flipped operator depends on the chosen basis,

\item the intermediate construction in Eq.~\eqref{eq:appendix_spin_flip_definition} is representation dependent,

\item the final concurrence \( C(\rho) \) is basis independent.

\end{itemize}

%--------------------------------------------------------
% D.7 - Manifestly Basis-Independent Formulations
%--------------------------------------------------------

\subsection*{Manifestly Basis-Independent Formulations}

For pure states, concurrence admits formulations that make the invariance more transparent.

Using the reduced density operator

\[
\rho_A
=
\mathrm{Tr_B}
\left(
|\psi\rangle\langle\psi|
\right),
\]

one obtains

\begin{equation}
C(|\psi\rangle)
=
2\sqrt{\det(\rho_A)}.
\label{eq:concurrence_determinant_appendix}
\end{equation}

Equivalently,

\begin{equation}
C(|\psi\rangle)
=
\sqrt{
2
\left(
1-\mathrm{Tr}(\rho_A^2)
\right)
}.
\label{eq:concurrence_purity_appendix}
\end{equation}

These expressions involve only invariant quantities such as determinants and traces, and therefore do not rely explicitly on any preferred basis representation.

\medskip

More fundamentally, concurrence for mixed states is defined through the convex-roof construction

\begin{equation}
C(\rho)
=
\min_{\{p_k,|\psi_k\rangle\}}
\sum_k
p_k
C(|\psi_k\rangle),
\label{eq:convex_roof_concurrence_appendix}
\end{equation}

where

\begin{equation}
\rho
=
\sum_k
p_k
|\psi_k\rangle\langle\psi_k|.
\label{eq:convex_decomposition_appendix}
\end{equation}

This formulation is intrinsically basis independent.

The spin-flip formalism of Eq.~\eqref{eq:appendix_spin_flip_definition} constitutes a special analytical construction that allows the explicit evaluation of the convex-roof optimization for two-qubit systems.

%--------------------------------------------------------
% D.8 - Physical Interpretation in Fiber-Based Quantum Channels
%--------------------------------------------------------

\subsection*{Physical Interpretation in Fiber-Based Quantum Channels}

The distinction between representation-dependent quantities and invariant quantities plays a central role in the modeling strategy developed throughout the present work.

\medskip

Local unitary transformations induced by fiber propagation,

\[
U_A \otimes U_B,
\]

may alter:

\begin{itemize}

\item local Bloch coordinates,

\item Stokes parameters,

\item tensor components,

\item fixed-basis Bell-state fidelities,

\item correlation functions evaluated along specific axes.

\end{itemize}

However, these transformations do not modify the intrinsic entanglement content of the bipartite state.

Consequently, the concurrence remains unchanged under coherent local polarization evolution.

\medskip

By contrast, non-unitary effects such as ensemble averaging, decoherence, or polarization-mode dispersion generate effective mixed states and therefore reduce concurrence.

This distinction establishes the conceptual transition

\[
\text{coherent local evolution}
\rightarrow
\text{effective decoherence}
\rightarrow
\text{entanglement degradation},
\]

which constitutes one of the central physical mechanisms investigated in the simulator developed in this work.
```
